\subsection{Первая программа}

Традиционно первая программа печатает приветствие. На дроне нет экрана, но есть консоль: в симуляторе это терминал редактора, в Pioneer Station --- окно «Консоль».

\begin{codefile}{lua/examples/02_sensors/getting_started_hello.lua}
\end{codefile}

Выберите язык \ui{Lua}, вставьте программу и нажмите \ui{Запустить}. Затем откройте терминал (значок \texttt{>\_} в заголовке редактора):

\begin{center}
\simshot[0.9\linewidth]{03-terminal.png}
\captionof{figure}{Терминал показывает всё, что программа напечатала через \texttt{print}, и сообщения автопилота.}
\end{center}

Что здесь произошло:
\begin{itemize}
    \item строки, начинающиеся с \texttt{--}, --- комментарии, Lua их не выполняет;
    \item \texttt{print(\dots)} печатает текст в консоль. Текст --- это \textbf{строка}, она записывается в кавычках;
    \item программа закончилась сразу после печати. Дрон не взлетел, потому что мы не дали ему ни одной команды.
\end{itemize}

\begin{taskbox}[Проверка связи]
Измените текст в кавычках на своё имя и запустите программу ещё раз. Затем добавьте вторую строку \texttt{print(2 + 2)} и убедитесь, что Lua умеет считать.
\end{taskbox}

\begin{infobox}[Где лежат примеры]
Каждая программа в книге подписана именем файла, например \texttt{lua/examples/02\_sensors/getting\_started\_hello.lua}. Это путь внутри папки \texttt{code/} курса. Те же файлы доступны в редакторе симулятора через список примеров, а текст программы можно скопировать прямо из PDF --- отступы и дефисы при копировании сохраняются.
\end{infobox}
